\uuid{3UMn}
\exo7id{552}
\titre{exo7 552}
\auteur{cousquer}
\organisation{exo7}
\datecreate{2003-10-01}
\isIndication{false}
\isCorrection{false}
\chapitre{Suite}
\sousChapitre{Suite définie par une relation de récurrence}
\cpge{mpsi-mp2i.nombres-reels-suites.c, pcsi.nombres-reels-suites.b, ptsi.nombres-reels-suites.b, mpsi-mp2i.nombres-reels-suites.d, pcsi.nombres-reels-suites.c, ptsi.nombres-reels-suites.c, mpsi-mp2i.nombres-reels-suites.h, pcsi.nombres-reels-suites.f, ptsi.nombres-reels-suites.f, mpsi-mp2i.nombres-reels-suites.i, pcsi.nombres-reels-suites.g, ptsi.nombres-reels-suites.g, mpsi-mp2i.limites-continuite.a, pcsi.limites-continuite.a, ptsi.limites-continuite.a, mpsi-mp2i.limites-continuite.b, pcsi.limites-continuite.b, ptsi.limites-continuite.b, mpsi-mp2i.derivabilite.a, pcsi.derivabilite.a, ptsi.derivabilite.a, bcpst1.suites-reelles-usuelles, bcpst1.suites-reelles, bcpst1.limites-continuite.a, bcpst1.limites-continuite.c, bcpst1.derivation.a}
\module{Analyse}
\niveau{L1}
\difficulte{}

\contenu{
\question{
Étudier suivant les valeurs données à $u_0$ appartenant à $\mathbb{C}$ les
suites~:
\begin{eqnarray*}
  u_{n+1} & = & \frac{u_n-2}{u_n +4}\\
  u_{n+1} & = & \frac{u_n+2}{u_n+1}\\
  u_{n+1} & = & \frac{-1}{u_n+1}
\end{eqnarray*}
}
}
